\documentclass[10pt,a4paper]{amsart}
\usepackage[margin=19mm]{geometry}
\usepackage{amsmath,amssymb,mathtools}
\usepackage[scr=rsfs,cal=euler]{mathalfa}
\usepackage{xfrac}
\usepackage[unicode,hidelinks,bookmarks=false]{hyperref}
\newcommand{\ie}{i.e.\ }
\newcommand{\cf}{cf.\ }
\newcommand{\resp}{respectively\ }
\newcommand{\Subsection}{\subsection}
\providecommand{\nobreakdash}{\nobreak\hbox{-}}
\setlength{\emergencystretch}{2em}
\multlinegap0pt
\usepackage{bm}
\usepackage{bbm}
\usepackage{dsfont}
\usepackage{xspace}
\usepackage{cancel}
\usepackage{mathtools}
\usepackage[shortlabels]{enumitem}
\usepackage{amsthm}
\makeatletter
\@ifundefined{lemma}{\newtheorem{lemma}{Lemma}}{}
\makeatother
\makeatletter
\newcommand{\di}[1]{\,\mathrm{d}#1}
\expandafter\ifx\csname proof\endcsname\relax
\newenvironment{proof}{\paragraph{Proof:}}{\hfill$\square$}
\fi
\makeatother
\title[Strongly Coupled Two-scale System with Nonlinear Dispersion:]{Strongly Coupled Two-scale System with Nonlinear Dispersion: Weak Solvability and Numerical Simulation}
\author{Vishnu Raveendran, Surendra Nepal, Rainey Lyons, Michael Eden, Adrian Muntean}
\date{Source: arXiv:2311.12251v4 (CC BY 4.0)}
\begin{document}
\maketitle

\noindent\emph{Source and licence.} Strongly Coupled Two-scale System with Nonlinear Dispersion: Weak Solvability and Numerical Simulation. Version arXiv:2311.12251v4 (CC BY 4.0), distributed under the Creative Commons Attribution 4.0 International licence, as recorded in the arXiv licence metadata for this exact version and shown on the abstract page https://arxiv.org/abs/2311.12251v4. Full rights notice: licenses/arxiv-2311.12251-CC-BY-4.0.txt.\par\smallskip

\noindent\textit{Editorial scope.} Counted unit: the Lemma whose source label is \texttt{L2} in the article (source0.tex lines 597--612, source label \texttt{L2}), delivered together with the article's own complete original proof of it (source0.tex lines 613--716, one \texttt{proof} environment of 104 lines). No mathematics is written, repaired or re-derived: statement and proof are the article's own text line for line. Required context retained uncounted: macrodiff (lines 331--336); A1 (lines 350--370); A3 (lines 353--356); L1 (lines 485--508); aux1 (lines 488--492); aux3 (lines 488--492); auxwf (lines 511--513); dbar (lines 575--577); L3 (lines 580--586). Every definition, display and label of those lines is retained unchanged. Omitted from this package and not redistributed: 1--330, 337--349, 357--484, 493--510, 514--574, 578--579, 587--596, 717--1328. Nothing inside a retained excerpt is omitted or altered: every retained body is delivered exactly as the article's lines are written, so no omission receipt of any kind accompanies this package. Citations: the retained excerpts carry no citation command at all, so no bibliography entry of the article is transcribed and no reference is redistributed. Reference adaptation: none was needed and none was made; every reference inside the delivered unit resolves inside it, so no unresolved reference or citation is produced. Printed slips kept literal: no printed word, symbol, hypothesis or number of the counted statement or of its proof was altered. Layout and class: the article's own class is \texttt{article} with packages including graphicx, wrapfig, lipsum, amsmath, amssymb, authblk, amsfonts, amsthm, geometry, etoolbox, enumerate, appendix, hyperref, fullpage; the wrapper is the lane's pinned \texttt{amsart} editorial wrapper at 10pt on A4, which restates the theorem-like environments the retained text needs and adds the article's own macro definitions that the retained text needs (\texttt{\textbackslash di}), with extra wrapper packages (bm, bbm, dsfont, xspace, cancel, mathtools, enumitem). The delivered numbering is the unit's own. Assets: no retained span contains a figure environment, a graphics-inclusion command, a PDF-page inclusion, a table or a TikZ picture; no image file is copied into this package, the package declares no asset dependency and no image file is redistributed with it. Licence: version arXiv:2311.12251v4 of this article is distributed under the Creative Commons Attribution 4.0 International licence, as recorded in the arXiv licence metadata for this exact version and shown on the abstract page https://arxiv.org/abs/2311.12251v4. A full rights notice is delivered at licenses/arxiv-2311.12251-CC-BY-4.0.txt. Characters: every retained excerpt is pure ASCII, so the delivered engine, XeLaTeX with its default Latin Modern text font, reports no missing character.
\begin{equation}\label{macrodiff}
     D^*(W):=\frac{1}{|Y|}\int_Y D(y)\left(I+\begin{bmatrix}
\frac{\partial w_1}{\partial y_1} & \frac{\partial w_2}{\partial y_1} \\
\frac{\partial w_1}{\partial y_2} & \frac{\partial w_2}{\partial y_2}
\end{bmatrix}\right)\di{y}
\end{equation}

\begin{enumerate}[label=({A}{{\arabic*}})]
  \item\label{A1} %(Ellipticity condition) 
The microscopic diffusion matrix satisfies $D\in (H^1_\#(Y)\cap L^\infty(Y))^{2\times2}$ and there exists $\theta>0$  such that 
 \begin{equation*}
 \theta |\eta |^{2} \leq
D\eta\cdot \eta \quad\text{for all}\  \eta\in \mathbb{R}^2 \; \mbox{and almost all } y\in Y;
 \end{equation*}
 
\item  $G_1,G_2:\mathbb{R}\rightarrow\mathbb{R}$ are locally Lipschitz functions, i.e., Lipschitz on compact sets;
 	 \label{A2}
  \item \label{A3} The microscopic drift velocity $B\in (H^1_\#(Y)\cap L^\infty(Y))^2$ satisfies
\begin{equation*}
    \begin{cases}
 	      \mathrm{div}_yB=0\hspace{.4cm}\mbox{in} \hspace{.4cm}   Y,\\
 	      B\cdot n_y =0 \hspace{.4cm}\mbox{on} \hspace{.4cm}  \Gamma_N;
 	\end{cases}
\end{equation*}

\item The reaction rate satisfies $f\in C^{\alpha,\frac{\alpha}{2}}((0,T)\times\Omega)$ and the initial condition $g\in C^{2+\alpha}(\Omega)$, for some $0<\alpha<1.$
\label{A4}
 \end{enumerate}

\begin{lemma}\label{L1}
Let $p \in \mathbb{R}$. Consider the following auxiliary problem: Find $W(p,\cdot)=(w_1(p,\cdot),w_2(p,\cdot)) \in \mathcal{W}^2$ satisfying
\begin{subequations}\label{system:cellproblem_with_p}
\begin{align}
    \mathrm{div}_y\left(- D \nabla_y w_i+pB w_i\right)&=\mathrm{div}_y(D e_i) &\mbox{in}& \hspace{.2cm}  Y,\label{aux1}\\
    \left( -D\nabla_y w_i+pB w_i\right)\cdot n_y &= \left( De_i\right)\cdot n_y &\mbox{on}& \hspace{.2cm} \Gamma_N,\label{aux2}\\
    w_i\, &\mbox{ is $Y$--periodic},&&\label{aux3}
    \end{align}
\end{subequations}
where $i\in\{1,2\}$.
{Assume \ref{A1}--\ref{A3}} holds, then 
\begin{enumerate}[label=(\roman*)]
    \item there exists a unique weak solution $w_i(p,\cdot)\in \mathcal{W}$   to the Problem \eqref{aux1}--\eqref{aux3},
 \item there exists a constant $C>0$ independent of $p$ such that
\begin{equation}
   \|\nabla_y w_i(p,\cdot)\|_{L^2(Y)}\leq C,\label{L1e1}
\end{equation}
\item there exists a constant $C>0$ such that, for all $p_1,p_2\in \mathbb{R}$, 
\begin{equation}\label{L1e2}
    \int_Y |\nabla_y( w_i(p_1,y)-w_i(p_2,y))|^2 \di{y}\leq C |p_1-p_2|^2.
\end{equation}
\end{enumerate}

\end{lemma}

\begin{equation}
      \int_{Y} \big(D\nabla_y w_i- pBw_i\big)\cdot \nabla_y \psi\di{y}=\int_{Y} \mathrm{div}_y( D e_i)\psi \di{y}-\int_{\Gamma_N} De_i\cdot n_y\psi \di{\sigma}\quad \label{auxwf}
\end{equation}

\begin{align}
    [\overline{D}(p)]_{i,j}:=\dfrac{1}{|Y|}\int_Y D(y)\left(e_j+\nabla_y w_j (p,y)\right) \cdot e_i\di{y},\label{dbar}
\end{align}

\begin{lemma}\label{L3}
    {Assume \ref{A1}--\ref{A3}} hold.
    There exists a constant $C>0$, independent of $p$ and $q$, such that 
 \begin{equation*}
     |[\overline{D}(p)]_{i,j}-[\overline{D}(q)]_{i,j}|\leq C |p-q|.
 \end{equation*}
\end{lemma}

\begin{lemma}\label{L2}
{Assume \ref{A1}--\ref{A3}} hold and let $M>0$.
For $p=(p_1,p_2)\in [-M,M]^2$, let $W_p:=(w_1(p_1,\cdot ),w_2(p_2,\cdot ))$ be the corresponding solutions of the auxiliary problem \eqref{aux1}--\eqref{aux3}. Then the macroscopic dispersion tensor $D^*(W_p)$ satisfies the following properties:
\begin{enumerate}[label=(\roman*)]
    \item There exists $\theta_M>0$ independent of $p$ such that 
         \begin{equation}\label{L2e1}
             \theta_M|\eta |^{2} \leq
             D^*(W_p)\eta \cdot \eta\quad \text{for all}\ \eta\in \mathbb{R}^2;
          \end{equation}

    \item There exist a constant $C>0$ independent of $p$ such that
         \begin{equation}\label{L2e2}
            | [D^*(W_p)]_{i,j}|\leq C,\quad (i,j\in\{1,2\}).
         \end{equation}
\end{enumerate}
\end{lemma}

\begin{proof}

    \par \textbf{\textit{(i)}} From the definition of $\overline{D}(\cdot)$ \eqref{dbar} and the definition of $D^*(\cdot)$ \eqref{macrodiff}, we have
    \begin{equation}
       [ D^*(W_p)]_{i,j}=[\overline{D}(p_i)]_{i,j},\label{Q}
    \end{equation}
    and hence,
\begin{equation}\label{xL2e1}
    [D^*(W_p)]_{i,j}=\frac{1}{|Y|}\int_Y D(y)\left(e_j+\nabla_y w_j(p_i,y)\right)\cdot e_i \di{y}.
\end{equation}
We consider the weak formulation for the Problem \eqref{aux1}--\eqref{aux3} (i.e., Equation \eqref{auxwf}) for index $j$, where we choose the test function $\psi=w_i$ and divide by $|Y|$:
%
\begin{equation}\label{L2e5}
    \frac{1}{|Y|}\int_{Y}\big(D\nabla_y w_j- p_jBw_j\big)\cdot \nabla_y w_i\di{y}
    =\frac{1}{|Y|}\int_Y\text{div}_y(D e_j)w_i\di{y}-\frac{1}{|Y|}\int_{\Gamma_N}\left( De_j\right)\cdot n_y w_i\di{\sigma}.
\end{equation}
Using integration by parts on $\int_Y\text{div}_y(D e_j)w_i\di{y}$ and noting the periodicity of $D(\cdot)$ and $w_i(\cdot)$, we arrive at
\begin{equation}\label{L2e6}
    \frac{1}{|Y|}\int_{Y}\big( D\nabla_y w_j- p_jBw_j+D e_j\big)\cdot \nabla_yw_i \di{y} =0.
\end{equation}
Adding \eqref{L2e6} to \eqref{xL2e1}, we see that
\begin{equation}\label{L2e7}
    [D^*(W_p)]_{i,j}=[A(p)]_{i,j}+[J(p)]_{i,j},
\end{equation}
where 
\begin{align}
    [A(p)]_{i,j}&:=\frac{1}{|Y|}\int_Y D(e_j+\nabla_yw_j)\cdot 
    (e_i+\nabla_y w_i)\di{y},\label{L2e8}\\
   \text{and } [J(p)]_{i,j}&:=-\frac{1}{|Y|}\int_{Y}p_jBw_j\cdot \nabla_y w_i\di{y}.\label{L2e9}
\end{align}
Using \ref{A3} and integration by parts on $\int_{Y}p_jBw_j\cdot \nabla_y w_i\di{y}$, we obtain
\begin{align}
    -[J(p)]_{i,j}&=-\frac{p_j}{|Y|}\int_{Y}\nabla_y(Bw_j)  w_i\di{y}+\frac{p_j}{|Y|}\int_{\partial Y}B\cdot n_y w_iw_j\di{\sigma}\nonumber\\
    &=-\frac{1}{|Y|}\int_{Y}p_jBw_i\cdot \nabla_y w_j\di{y}\nonumber\\
    &=[J(p)]_{j,i}.\label{L2e10}
\end{align}
Combining \eqref{L2e8} and \eqref{L2e10} shows that $A(p)$ and $J(p)$ are the symmetric and skew-symmetric parts of $D^*(W_p)$.
Since $J(p)$ is a $2\times 2$ skew-symmetric matrix, we have 
\begin{equation}\label{L2e11}
    J(p) \eta\cdot \eta=0\quad\text{for all}\ \eta\in \mathbb{R}^2.
\end{equation}
So, to prove \eqref{L2e1} it is enough to establish that $A(p)$ is in fact positive definite {for all $p\in [-M,M]^2$}. {For that by the standard arguments we first show that for a fixed $p$,  $A(p)$ is positive definite. }
Since $A(p)$ is symmetric, using the structure \eqref{L2e8} we have for all $\eta=(\eta_1,\eta_2)\in \mathbb{R}^2$, 
\begin{align}
A(p)\eta\cdot\eta&=\eta_1^2\frac{1}{|Y|}\int_Y D(e_1+\nabla_yw_1)\cdot 
    (e_1+\nabla_y w_1)\di{y}\nonumber\\
    &\hspace{1cm}+2\eta_1 \eta_2\frac{1}{|Y|}\int_Y D(e_1+\nabla_yw_1)\cdot 
    (e_2+\nabla_y w_2)\di{y}\nonumber\\
    &\hspace{1.5cm}+\eta_2^2 \frac{1}{|Y|}\int_Y D(e_2+\nabla_yw_2)\cdot 
    (e_2+\nabla_y w_2)\di{y}.\nonumber
    \end{align}
    Distributing $\eta_1$ and $\eta_2$ inside the integrals, we get
    \begin{align}
    A(p)\eta\cdot\eta&=\frac{1}{|Y|}\int_Y D(\begin{bmatrix} \eta_1 \\ 0  \end{bmatrix}+\eta_1\nabla_yw_1)\cdot 
    (\begin{bmatrix} \eta_1 \\ 0  \end{bmatrix}+\eta_1\nabla_y w_1)\di{y}\nonumber\\
    &\hspace{1cm}+2\frac{1}{|Y|}\int_Y D(\begin{bmatrix} \eta_1 \\ 0  \end{bmatrix}+\eta_1\nabla_yw_1)\cdot 
    (\begin{bmatrix} 0 \\ \eta_2  \end{bmatrix}+\eta_2\nabla_y w_2)\di{y}\nonumber\\
    &\hspace{1.5cm}+ \frac{1}{|Y|}\int_Y D(\begin{bmatrix} 0 \\ \eta_2  \end{bmatrix}+\eta_2\nabla_yw_2)\cdot 
    (\begin{bmatrix} 0 \\ \eta_2  \end{bmatrix}+\eta_2\nabla_y w_2)\di{y}\label{xL2e2}.
      \end{align}
      Rearranging the right side of the identity \eqref{xL2e2} and using \ref{A1} yields
    \begin{align}
           A(p)\eta\cdot\eta&=\frac{1}{|Y|}\int_Y D(\begin{bmatrix} \eta_1 \\ \eta_2  \end{bmatrix}+\sum_{i=1}^2\eta_i\nabla_yw_i)\cdot 
    (\begin{bmatrix} \eta_1 \\ \eta_2  \end{bmatrix}+\sum_{i=1}^2\eta_i\nabla_yw_i)\di{y}\nonumber\\
    &\geq \theta\frac{1}{|Y|}\int_Y |\eta+ \sum_{i=1}^2\eta_i\nabla_yw_i)|^2\;\di{y}\nonumber\\
    &\geq 0.\label{L2e12}
\end{align}
If $A(p)\eta\cdot\eta=0$ for some $\eta\in \mathbb{R}^2$, then we have
\begin{align*}
|\eta+ \sum_{i=1}^2\eta_i\nabla_yw_i)|^2=0 \,\,\,\mbox{for almost every }\, y\in Y.
\end{align*}
Consequently, this yields
\begin{equation}
\sum_{i=1}^2\eta_iw_i=C-\eta\cdot y,\label{L2e13}
\end{equation}
for some constant $C $.
Since the $w_i$ are periodic, identity \eqref{L2e13} is only possible for $\eta=0$.
Hence, $A(p)\eta\cdot \eta=0$ only for $\eta= 0$.
Therefore, combining this with \eqref{L2e11}, we get
\begin{equation}\label{L2e15}
    D^*(W_p)\eta\cdot \eta> 0.
\end{equation}
for all $p\in [-M,M]^2$ and all $\eta\neq0$.
 Define the function $F:[-M,M]^2\times \partial B(0,1)\rightarrow \mathbb{R}$ as 
\begin{equation*}
    F(p,\xi):= D^*(W_p)\xi\cdot \xi \,\,\,\mbox{for all}\; \xi\in \partial B(0,1).
\end{equation*}
     We have from Lemma \ref{L3} that  the function $p_i\mapsto \overline{D}(p_i)$ is continuous.
     Using the definition \eqref{Q}, we see that the function $p\mapsto D^*(W_p)$ is also continuous.
     Therefore, $F$ is continuous and $F(p,\xi)>0$ over the compact set $ ([-M,M]^2\times \partial B(0,1))$. Hence, there exist a $\theta_M>0$ such that
\begin{equation*}
    F(p,\xi)>\theta_M
\end{equation*}
for all $(p,\xi)\in  ([-M,M]^2\times \partial B(0,1))$.
So, for any $\eta\neq 0$ we have 
\begin{align*}
    D^*(W_p)\frac{\eta}{|\eta|}\cdot \frac{\eta}{|\eta|}&>\theta_M,\\
     \text{i.e., } D^*(W_p)\eta\cdot \eta&>\theta_M|\eta|^2.
\end{align*}
\textbf{\textit{(ii)}} The inequality \eqref{L2e2} follows directly from \ref{A1}, $(ii)$ of Lemma \ref{L1}, and the identity 
\begin{equation*}
    [D^*(W_p)]_{i,j}=\frac{1}{|Y|}\int_Y D\left(e_j+\nabla_y w_j\right)\cdot e_i\di{y}
\end{equation*}together with the Cauchy--Schwartz inequality.
\end{proof}

\end{document}
